\section{What the Genetic Model Assumes and Predicts}
\label{sec:fisher-model}
\subsection{The Outcome and the Genetic Estimand}
\label{sec:model-components}
The benchmark is an additive model for a specified phenotype, not a general
theory of every outcome called social status. Write the standardized phenotype
as
\[
 P_i=hA_i+e_P E_i,\qquad e_P=\sqrt{1-h^2},
\]
where $A_i$ and $E_i$ have zero means, unit variances, and zero within-person
covariance. The additive genetic contribution is $hA_i$, with variance
$h^2$. Define $m=\operatorname{Corr}(A_M,A_F)$ and
$r=\operatorname{Corr}(P_M,P_F)$ for spouses. The sexes share component
variances, and those variances are stationary across generations.
The formulation follows \citet{Fisher1918} and the explicit treatment of \citet{OttoEtAl1994}. Appendix~\ref{app:fisher} gives the primitive assumptions and derivations;
there $\rho=r$ and $\rho_A=m$.

The intended causal genetic contribution must be distinguished from a
population predictor. In modern genomics, a direct effect originates in
the individual's own genotype; an indirect effect originates in another
person's genotype, apart from transmission of the individual's alleles.
The direct effect can include parental, teacher, or peer responses to the
individual's genetically influenced characteristics. It is not an effect
that holds all social responses fixed
\citep[Sections~3.1 and~3.4]{BenjaminEtAl2026}. This definition matters both
for the covariance exclusions and for Clark's later interpretation of
``social'' effects.

\subsection{The Strength of the Environmental Covariance Exclusions}
\label{sec:environmental-exclusions}

The zero-covariance assumptions do much of the substantive work. Within
a person, $\operatorname{Cov}(A_i,E_i)=0$ excludes systematic association
between the direct genetic contribution and the remaining environmental
contribution. Across biological relatives, the model excludes shared
environmental contributions: $\operatorname{Cov}(E_i,E_j)=0$. In the birth
process of Appendix~\ref{app:fisher}, each child's environment is also uncorrelated with parental
genes and environments and with siblings' genes. In particular, for a
parent $p$ and offspring $o$, and distinct siblings $1,2$,
\[
 \operatorname{Cov}(A_p,E_o)=\operatorname{Cov}(E_p,E_o)=0,
 \qquad
 \operatorname{Cov}(A_1,E_2)=\operatorname{Cov}(E_1,E_2)=0.
\]
These are substantive exclusions, not consequences of additive genetics
or of Mendelian segregation.

For education, occupation, earnings, and wealth, this is an exceptionally
restrictive starting point. It excludes any net contribution to familial
resemblance from shared schools, neighborhoods, resources, contacts, or
inherited wealth beyond the pathways already assigned to genetic effects.
It also excludes additional association between the child's inherited
endowment and the environmental advantages or disadvantages supplied by
the family. Individual environments may have large effects in this model,
but the residual environmental influences must be uncorrelated along the
specified family relationships. A large environmental variance therefore
does not make the model permissive about environmental transmission.

As \citet[pp.~10--11]{Goldberger1978} emphasizes, the asymmetry is striking:
premarital environments may resemble one another among spouses, while
shared rearing supplies no additional source of resemblance among
siblings or between parents and children. For otherwise comparable
biological relatives, being raised together rather than apart contributes
nothing to the predicted correlation. For Clark's social outcomes, these
restrictions demand unusually strong empirical justification. They cannot
be established by fitting equations that impose them from the outset.
Section~\ref{sec:genomic-evidence} explains why modern genomic evidence
already challenges the exclusion of environmental confounding.

The behavioral asymmetry deserves emphasis. Environmental contributions
to parents' phenotypes help determine whom they marry, and assortment
thereby associates those environments with the genes their children
inherit. Yet the model gives each child a fresh environmental draw,
disconnected from parental circumstances and from the sibling's
environmental draw. Parents' environments thus enter the process of mate
selection but contribute no additional resemblance between the children
they raise. This construction is mathematically coherent: correlation with
a parent's environment need not imply correlation with one's own when
the model disconnects the two. Its behavioral plausibility is a separate
question. For education, occupation, and wealth, excluding the family
processes that could connect parental circumstances to children's
opportunities requires evidence, not merely a consistent covariance matrix.
Goldberger's question is therefore substantive: what account of family
behavior supports these restrictions \citep[pp.~36--37]{Goldberger1978}?

Two qualifications are essential. First, it would be incorrect to say
that \emph{every} cross-person genetic--environmental covariance outside
marriage is zero. Phenotypic assortment generates correlation between one
spouse's environment and the other's genes; inheritance carries that
association to their child. The path runs from association with the
partner's genes to inheritance of those genes by the child; it is not a
causal effect of the parent's environment on the child's genes. Nor is it
transmission of the parent's environment to the child's environment.
Indeed, the derivation in Appendix~\ref{app:fisher} gives
$\operatorname{Cov}(E_M,A_o)=r h e_P/2$, even though
$\operatorname{Cov}(A_M,E_o)=0$. The restriction excludes environmental
transmission and additional covariance, while retaining the covariance
already implied by the mating mechanism.

Second, a direct genetic effect can include environmental responses to
the individual's genotype. Such mediation belongs to the direct effect
under the causal definition used in modern genomics; it is not automatically
an omitted environmental term. Conversely, defining the genetic component
as a population linear predictor can make its residual orthogonal by
construction, but does not make that predictor a direct causal genetic
contribution. Clark cannot obtain the causal interpretation merely by
choosing an orthogonal statistical decomposition. This distinction is
central to Sections~\ref{sec:genomic-evidence} and~\ref{sec:interpretation}.


\subsection{The Mating Mechanism Restricts the Parameters}
Under primary phenotypic assortment, spouses match on $P$ rather than
separately on its genetic and environmental components. With the conditional
linearity and independence assumptions stated in Appendix~\ref{app:fisher},
\begin{equation}
 m=h^2r.\label{eq:overview-restriction}
\end{equation}
The relation is derived from the matching mechanism; it is not an optional
additional assumption or a consequence that requires an intergenerational
equilibrium argument. A high phenotype predicts a genetic value of $hP$;
applying that prediction on both sides of a marriage gives the factor $h^2$.

Assortment on genetic value instead gives $r=h^2m$ under independent
residual environments. Assortment on several traits can generate still
other restrictions. Clark's appeal to matching on ``something deeper''
therefore requires a specification \citep[pp.~16,75]{Clark2026}: either it is the same phenotypic
model with a measurement equation, or it changes the mating process and
requires new kinship formulas. These possibilities cannot be exchanged
while retaining the original parameter interpretation.

\subsection{Inheritance and the Predictions for Relatives}
The child receives half each parent's additive genetic value plus a
mean-zero segregation deviation:
\[
 A_o=\tfrac12(A_M+A_F)+S_o.
\]
Independent segregation across births and stationary unit genetic variance
require $\operatorname{Var}(S_o)=(1-m)/2$. The child's environmental draw
is fresh in the sense specified above. Under phenotypic assortment the
resulting correlations are
\[
 \operatorname{Corr}(P_p,P_o)=\frac{h^2(1+r)}2,
 \qquad
 \operatorname{Corr}(P_1,P_2)=\frac{h^2(1+m)}2
\]
for a parent and child and for full siblings, respectively. Thus
\begin{equation}
 \operatorname{Corr}(P_p,P_o)-\operatorname{Corr}(P_1,P_2)
 =\frac{h^2r(1-h^2)}2>0
 \label{eq:overview-gap}
\end{equation}
when assortment is positive and $0<h^2<1$. More distant moments require
additional assumptions about incoming spouses and shared ancestors,
not merely Mendelian transmission. The complete derivation exposes where
each exclusion enters. The next question is whether these latent
predictions survive the way Clark's outcomes are recorded and fitted.
